\chapter{把一个服务放到公网}
\label{ch:public-service}
\chapterlead{对外提供服务要解决四件事：进程如何常驻、域名如何指向、证书如何续期、代码如何更新。本章用一个静态站点与一个后端服务走完整条链路。}

\section{起步：让一个进程常驻}

先用 systemd 跑一个最小的后端，确认它能在重启后自动恢复。

\begin{nixcode}
services.nginx.enable = true;
services.nginx.virtualHosts."example.com" = {
  forceSSL = true;
  enableACME = true;
  root = "/srv/www/example";
};

security.acme = {
  acceptTerms = true;
  defaults.email = "weiwen@example.com";
};
\end{nixcode}

重建后，nginx 自动申请证书、配置 80 到 443 的跳转，并由 systemd 定时续期。这一步把“域名、证书、静态文件”三件事一起解决。

\begin{shellsession}
$ curl -I https://example.com
$ systemctl status nginx
$ journalctl -u acme-example.com -n 20
\end{shellsession}

\section{反向代理一个后端}

后端程序监听高位端口，由 nginx 转发。

\begin{nixcode}
services.nginx.virtualHosts."api.example.com" = {
  forceSSL = true;
  enableACME = true;
  locations."/" = {
    proxyPass = "http://127.0.0.1:3000";
    proxyWebsockets = true;
  };
};

systemd.services.api = {
  description = "API service";
  wantedBy = [ "multi-user.target" ];
  after = [ "network.target" ];
  serviceConfig = {
    ExecStart = "/srv/api/bin/api";
    User = "deploy";
    WorkingDirectory = "/srv/api";
    Restart = "on-failure";
    RestartSec = 3;
    StateDirectory = "api";
  };
};
\end{nixcode}

\texttt{User} 让服务以部署账户运行，\texttt{StateDirectory} 自动创建 \texttt{/var/lib/api} 并设好权限。\texttt{Restart = "on-failure"} 让崩溃的进程自动拉起，\texttt{RestartSec} 避免立即重启导致的循环。

\begin{shellsession}
$ sudo systemctl restart api
$ curl -s localhost:3000/health
$ journalctl -u api -f
\end{shellsession}

\section{部署流程}

三条常见路径，按复杂度排列。

\begin{booktable}{\textwidth}{J{0.815} J{1.185}}
\toprule
方式 & 适用 \\
\midrule
\endhead
推送到服务器后重建 & 单机、单人维护，配置即部署 \\
构建产物由 CI 上传 & 需要编译或测试，部署只做替换 \\
容器镜像 & 依赖复杂，需要与宿主隔离 \\
\bottomrule
\end{booktable}

第一种最简单：服务器上 \texttt{git pull} 后执行 \texttt{nixos-rebuild switch}。适用于配置改动；程序本身若需要编译，应在本地或 CI 完成，服务器只接收产物。

第二种把构建与部署分开。CI 产出 store 路径或压缩包，服务器拉取后切换。好处是服务器不需要完整的构建环境。

\begin{shellsession}
$ nix build .#api
$ nix copy --to ssh://srv1 ./result
$ ssh srv1 'sudo /run/current-system/bin/switch-to-configuration switch'
\end{shellsession}

\section{上线前的检查清单}

\begin{enumerate}
  \item 服务以非 root 用户运行，且只监听必要地址。
  \item 证书能自动续期，续期任务有监控。
  \item 日志写入 journald，能被 \texttt{journalctl -u} 查到。
  \item 重启后服务自动恢复，重建后配置不丢。
  \item 有一次成功的回滚演练，知道回到哪一代。
  \item 数据目录在备份范围内，见第\ref{ch:backup}章。
\end{enumerate}

\begin{notebox}[不要跳过回滚演练]
上线前切换一次再回滚一次，确认两条路都通。真正出故障时，没有演练过的回滚流程通常会在细节上卡住：服务名记错、旧代的编号找不到、依赖的目录权限不对。
\end{notebox}

\section{练习}

\begin{enumerate}
  \item 用 nginx 与 ACME 部署一个静态站点，确认证书自动签发。
  \item 为一个本地后端配置反向代理，验证 WebSocket 能正常工作。
  \item 写一个 systemd 服务，以非 root 用户运行，重启后自动恢复。
  \item 用 \texttt{nix copy} 把构建产物传到服务器，切换一次并回滚一次。
  \item 按六项检查清单逐条核对上面的部署，记录每一条的证据。
\end{enumerate}

\begin{checkpoint}
本章结束时，应能：
\begin{itemize}
  \item 用 nginx 与 ACME 部署静态站点，说明证书续期机制；
  \item 配置反向代理并处理 WebSocket；
  \item 写出带用户、工作目录、自动重启的服务定义；
  \item 在三种部署方式中选择一种，说明取舍；
  \item 用检查清单验收一次上线，并完成回滚演练。
\end{itemize}
\end{checkpoint}
